%==============================================================================
% USENIX Security submission scaffold.
%
% NOTE ON DOCUMENT CLASS: For the actual submission, substitute the official
% USENIX class `usenix2019_v3` (\documentclass[letterpaper,twocolumn,10pt]{usenix2019_v3}).
% We use the stock `article` two-column layout here so the scaffold compiles
% with a plain pdflatex install (no external .cls, no external .bib).
%
% INTEGRITY NOTE: This file is a SCAFFOLD. Every result number is a \todo{}
% placeholder naming the experiment that must fill it from REAL runs. Do NOT
% populate any table by hand. See the \todo list summary in the build log.
%==============================================================================
\documentclass[letterpaper,twocolumn,10pt]{article}

\usepackage{booktabs}
\usepackage{graphicx}
\usepackage{xcolor}
\usepackage{url}
\usepackage{amsmath}
\usepackage{amssymb}
\usepackage[colorlinks=true,allcolors=blue]{hyperref}

% ----------------------------------------------------------------------------
% Placeholder macro. Every empirical number that must come from a REAL run is
% wrapped so it is impossible to miss and impossible to mistake for a result.
% ----------------------------------------------------------------------------
\newcommand{\todo}[1]{\textcolor{red}{[TODO: #1]}}

% Compact spacing helpers for margins in this scaffold.
\setlength{\parskip}{0pt}

\title{\vspace{-3ex}Forcing vs.\ Memorization:\\
A Negative-Control Audit of Optimization-Based\\ PII Extraction from Language Models}

\author{Anonymous Submission\\ \small Under review at USENIX Security}
\date{}

\begin{document}
\maketitle

%==============================================================================
\begin{abstract}
%==============================================================================
Gradient-based prompt optimization (e.g., Greedy Coordinate Gradient, GCG) is
increasingly proposed as a way to audit what personally identifiable
information (PII) a language model has memorized: if an optimizer can drive the
model to emit a known record, the record is deemed extractable. Using a
controlled corpus with complete ground truth and \emph{negative-control}
individuals who never appear in training, we show that this reasoning is unsafe.
Optimization-based extraction conflates \emph{memorization recall} with the
optimizer's ability to \emph{force} arbitrary outputs: a naive GCG attack
elicits PII from never-trained individuals at nearly the same rate as from
trained ones (\todo{neg-ctrl vs.\ real EMR, e.g. 44\% vs.\ 45.5\% on GPT-2 --- Table~\ref{tab:main}}),
and its success is \emph{invariant} to how often a record occurred in
training---both hallmarks of forcing, not recall. An unconstrained white-box
soft-prompt attack forces \todo{soft-prompt EMR, ${\approx}$100\%} of targets.

We argue that a memorization audit is only meaningful \emph{net of a negative
control}, and make that control a first-class instrument. We report the
\textbf{negative-control-adjusted} extraction rate ($\text{optimized} -
\text{negctrl}$) as the true memorization signal; place the attack spectrum
(fixed prompts, PII-Scope/PII-Compass discovery, GCG, soft-prompt) on a
\emph{forcing axis} scored by both raw and adjusted rates; and show that a
fluency-constrained attack lowers the negative-control rate, trading raw
extraction for validity. Across model scales we
\todo{calibrate the adjusted signal / quantify how much raw rates overstate memorization --- Table~\ref{tab:main}}.
Our results reframe optimization-based privacy auditing around a control that
most prior extraction work omits.

\end{abstract}

%==============================================================================
\section{Introduction}
\label{sec:intro}
%==============================================================================

Organizations increasingly fine-tune language models (LMs) on corpora derived
from customer records, employee databases, and communication archives.
Personally identifiable information (PII) absorbed during training persists in
model weights and can be elicited long after ingestion. Verbatim email
addresses have been extracted from GPT-2~\cite{carlini2021extracting},
production systems reproduce training data at
scale~\cite{nasr2023scalable}, and structured PII is
specifically at risk~\cite{lukas2023analyzing}. The operational question for a
data owner is therefore not \emph{whether} leakage is possible in the
abstract, but \emph{how much} of a specific corpus a realistic---or a
worst-case---adversary can pull back out.

\paragraph{Optimization audits and the forcing problem.}
A natural way to audit this risk is to \emph{optimize} for it: given a record
the auditor already holds, run a gradient-based prompt attack (GCG and its
white-box relatives) that searches for an input driving the model to emit the
record verbatim. If the search succeeds, the record is called extractable. The
premise is that a successful optimization reveals \emph{memorization}. We show
this premise is unsafe. Gradient optimization over a large discrete (or
continuous) input space is powerful enough to \emph{force} a small model to emit
a target string \emph{whether or not the model ever saw it}. On a controlled
corpus with known ground truth, a naive GCG attack elicits PII from
\emph{negative-control} individuals---people who never appear in
training---at almost the same rate as from trained individuals
(\todo{44\% vs.\ 45.5\% on GPT-2}), and its success does \emph{not} increase
with how often a record occurred. These are signatures of forcing, not recall:
the ``extraction rate'' partly measures the \emph{attack's} expressivity rather
than the \emph{model's} memory.

\paragraph{The confirmation setting, and why the control is essential.}
We work in a \emph{confirmation} (auditing) setting: the auditor holds the
ground-truth record $t$ and has white-box access, and asks how strongly $M$ can
be made to reproduce $t$. This is the setting an organization auditing its own
corpus occupies, and---correctly measured---it yields a worst-case upper bound
on elicitability. But an unconstrained optimizer inflates that ``upper bound''
with forcing. The remedy is a control that most extraction work omits: run the
identical attack against records that were \emph{never trained on}, and subtract
their success. What survives---the \textbf{negative-control-adjusted} rate---is
the extraction attributable to memorization rather than to the optimizer.

\paragraph{Contributions.}
\begin{enumerate}
\item \textbf{Forcing vs.\ memorization.} We show optimization-based PII
extraction conflates the two, and give the diagnostics that separate them---%
negative controls and a frequency-response test on a controlled corpus with
complete ground truth (Sections~\ref{sec:method},~\ref{sec:results}).
\item \textbf{A corrected metric.} We define and report the
negative-control-adjusted extraction rate ($\text{optimized} -
\text{negctrl}$) as the memorization-attributable signal, and show that raw
optimization-extraction rates overstate memorization.
\item \textbf{An attack spectrum on a forcing axis.} We place fixed prompts,
PII-Scope~\cite{piiscope2024} / PII-Compass~\cite{piicompass2024} discovery,
GCG, and a white-box soft-prompt attack on a single axis scored by \emph{both}
raw and adjusted rates, revealing that the most powerful attacks force the most
(the soft-prompt attack forces \todo{${\approx}$100\%}).
\item \textbf{Constraint reduces forcing.} A fluency-regularized (adaptive)
GCG lowers the negative-control rate, trading raw extraction for validity, and
connects the forcing story to a perplexity-filter defense
(Section~\ref{sec:defense}).
\end{enumerate}

We emphasize what this paper is \emph{not}: it is not a new, stronger attack for
pulling out more PII. Its contribution is methodological---how to audit
memorization with optimization attacks \emph{honestly}, by controlling for the
forcing they introduce.

%==============================================================================
\section{Threat Model \& Problem Setup}
\label{sec:threat}
%==============================================================================

\subsection{Auditor Knowledge}
We frame the work as data auditing. An organization holds a training corpus
$D$ containing PII records and a model $M$ with parameters $\theta$ fine-tuned
on $D$. The auditor has (i) \textbf{white-box access} to $M$ (realistic for
open-weight targets and for an organization auditing its own model before
deployment) and (ii) \textbf{ground-truth knowledge of the target records}
$t \in D$ (the organization controls the corpus). Assumption (ii) is what
makes this a \emph{confirmation} setting and is stated here in the threat
model, not buried in an appendix.

\subsection{The Three-Point Auditing Spectrum}
We situate three attacker capabilities on a single axis of strength
(Figure~\ref{fig:spectrum}):
\begin{itemize}
\item \textbf{Fixed prompts (lower bound).} Handcrafted direct, completion,
few-shot, and template prompts. This is current audit practice and
under-reports risk.
\item \textbf{Multi-query / few-shot discovery (realistic middle).}
PII-Scope-style~\cite{piiscope2024} multi-query and soft-prompt attacks and
PII-Compass-style~\cite{piicompass2024} grounding with known prefixes. These
approximate a resourceful but non-gradient adversary.
\item \textbf{Gradient-optimized confirmation (upper bound).} An adapted GCG
attack that optimizes a prompt suffix directly against the known target. This
bounds what any prompt-based method can elicit.
\end{itemize}

\begin{figure}[t]
\centering
\fbox{\parbox{0.92\linewidth}{\centering \vspace{4ex}
\todo{Figure 1 --- threat-model spectrum diagram. Horizontal axis of attacker
strength: fixed prompts (lower bound) $\rightarrow$ multi-query/few-shot
discovery, PII-Scope/PII-Compass (realistic middle) $\rightarrow$
gradient-optimized confirmation (upper bound). Annotate auditor knowledge
(holds target $t$, white-box $M$) and mark the ``auditing gap'' as the span
from fixed-prompt report to upper bound.}
\vspace{4ex}}}
\caption{The PII auditing spectrum. The auditor holds the ground-truth record
and measures its elicitability; the gap between fixed-prompt audits and the
gradient-optimized upper bound is what we calibrate.}
\label{fig:spectrum}
\end{figure}

\subsection{The Confirmation Attack, Formalized}
Given model $M$ with parameters $\theta$ and a known target PII sequence
$t = (t_1, \ldots, t_n)$ drawn from the training corpus, the auditor seeks an
optimized prompt $p^\star$ over a length-$k$ token budget that maximizes the
model's probability of reproducing $t$:
\begin{equation}
p^\star = \arg\max_{p \in V^k} \, \log P(t \mid p, M),
\label{eq:objective}
\end{equation}
where $V$ is the vocabulary. Because $t$ is \emph{known}, this is a
well-posed optimization; the target is a fixed string, not a quantity to be
discovered. Equivalently we minimize the negative log-likelihood
\begin{equation}
\mathcal{L}(p) = -\sum_{i=1}^{n} \log P\!\left(t_i \mid p, t_1, \ldots, t_{i-1}; M\right).
\label{eq:loss}
\end{equation}
A record is \emph{confirmed extractable} if the model's greedy decoding under
$p^\star$ reproduces $t$ under the unified exact-match rule of
Section~\ref{sec:method}.

\subsection{The Forcing Problem and the Adjusted Bound}
\label{sec:forcing}
Naively, the fraction of confirmed records is an ``auditing upper bound.'' But
Eq.~\ref{eq:objective} is optimized over a large input space $V^k$, and for a
sufficiently expressive attack the maximizer $p^\star$ can drive $M$ to emit
$t$ even when $t$ was \emph{never in the training corpus}: the success then
reflects the attack's expressivity---\emph{forcing}---not the model's memory.
Confirmed-extractability therefore upper-bounds memorization only \emph{net of
forcing}.

We estimate the forcing floor directly. Let $D$ be the trained records and let
$C$ be a set of \emph{negative-control} records, drawn from the same
distribution but \emph{never} included in training. Running the identical attack
on both, let $\mathrm{EMR}(D)$ and $\mathrm{EMR}(C)$ be the confirmed-extraction
rates. Because control records carry no memorization signal, $\mathrm{EMR}(C)$
is a direct estimate of the attack's forcing rate. We define the
\textbf{negative-control-adjusted} rate
\begin{equation}
\mathrm{EMR}_{\text{adj}} \;=\; \mathrm{EMR}(D) - \mathrm{EMR}(C),
\label{eq:adjusted}
\end{equation}
and treat $\mathrm{EMR}_{\text{adj}}$---not $\mathrm{EMR}(D)$---as the
memorization-attributable extraction and hence the defensible auditing signal. A
second, orthogonal check is the \emph{frequency response}: genuine recall should
rise with how often a record occurred, whereas forcing is frequency-invariant
(Section~\ref{sec:results}). An attack whose $\mathrm{EMR}(C)\!\approx\!
\mathrm{EMR}(D)$ and whose success is flat in frequency is measuring forcing, and
its raw rate must not be reported as memorization.

%==============================================================================
\section{Related Work}
\label{sec:related}
%==============================================================================

\subsection{Memorization and Extraction}
LMs memorize training data, sometimes necessarily for long-tailed
generalization~\cite{feldman2020neural}, with memorization rising in model
scale, duplication, and example
frequency~\cite{tirumala2022memorization, carlini2022quantifying}.
Carlini et al.~\cite{carlini2021extracting} first extracted training data from
GPT-2; \cite{carlini2023quantifying} quantified the duplication-extractability
relationship; Nasr et al.~\cite{nasr2023scalable} scaled extraction to
production systems; and Biderman et al.~\cite{biderman2023emergent} studied
emergent and predictable memorization. Lukas et
al.~\cite{lukas2023analyzing} examined PII leakage specifically. This line
establishes that leakage \emph{occurs}, but reports raw extraction rates and
rarely runs the identical attack against \emph{never-trained} records; the
forcing contribution to those rates therefore goes unmeasured---the gap we
close with an explicit negative control (Section~\ref{sec:forcing}).

\subsection{Adversarial Prompt Optimization}
Zou et al.~\cite{zou2023universal} introduced GCG, gradient-guided discrete
optimization of an adversarial suffix. Follow-on work made such attacks fluent
or stealthy---AutoDAN~\cite{liu2023autodan} and
COLD-Attack~\cite{guo2024cold}---which is directly relevant to our adaptive
adversary since fluency regularization is exactly what defeats a perplexity
filter. Schwarzschild et al.~\cite{schwarzschild2024rethinking} recast
memorization as compression via the Adversarial Compression Ratio (ACR), also
GCG-based but on generic text with no PII, no discovery comparison, and no
defense. Common to the whole line is that the very expressivity that makes GCG a
``universal'' jailbreak---optimizing over $V^k$ to elicit an arbitrary
output~\cite{zou2023universal}---also lets it \emph{force} outputs a model never
memorized. Whether a confirmed extraction reflects recall or forcing cannot be
read off the success rate alone; it requires a never-trained control, which we
supply.

\subsection{The 2024--25 PII-Attack Line}
Three recent works define the space our contribution must be differentiated
from, and which the prior version of this paper failed to cite.
\textbf{PII-Scope}~\cite{piiscope2024} benchmarks PII extraction attacks
including a white-box soft-prompt optimization and multi-query variants,
reporting large multi-query gains. \textbf{PII-Compass}~\cite{piicompass2024}
improves targeted extraction by grounding prompts with true-prefix context (a
known-target confirmation-style method). \textbf{$(n,p)$-discoverable
extraction}~\cite{hayes2025np} formalizes a probabilistic worst-case notion of
extractability under $n$ queries at success probability $p$, giving the formal
anchor for a worst-case audit.

\subsection{Delta vs.\ Prior Work}
Our contribution is methodological---a control and a corrected metric, not a new
attack. \textbf{The primary delta} is that we run every attack against a
never-trained negative control and report the adjusted rate
(Eq.~\ref{eq:adjusted}); to our knowledge no prior PII-extraction study does
this, and it changes the conclusions (a raw rate we would otherwise report as
memorization is largely forcing). We are further differentiated on the following
axes. \textbf{vs.\ ACR}~\cite{schwarzschild2024rethinking}:
ACR is generic-text with no PII structure, no attack-spectrum comparison, and
no defense; we add PII-specific field-type/frequency structure and the
calibrated gap across the spectrum. \textbf{vs.\ PII-Scope /
PII-Compass}~\cite{piiscope2024,piicompass2024}: they provide attacks but no
gradient \emph{upper bound}, no formal $(n,p)$ anchoring, and no adaptive
defense loop; we benchmark \emph{directly against them} on shared models
(Section~\ref{sec:results}) so the gap is measured, not asserted.
\textbf{vs.\ the jailbreak adaptive line}~\cite{liu2023autodan,guo2024cold}:
all jailbreak-only; our PII-specific finding is the low-perplexity-target vs.\
low-perplexity-prompt tension (Section~\ref{sec:defense}).
\textbf{vs.\ known predictors}: we explicitly demote perplexity, compression,
and entity density to \emph{confirmatory}
(\cite{carlini2023quantifying, biderman2023emergent, schwarzschild2024rethinking})
and claim as new only field-type/structure-as-extractability and
type-token-ratio erosion under optimization.

%==============================================================================
\section{Methodology}
\label{sec:method}
%==============================================================================

\subsection{Confirmation via Adapted GCG}
We adapt GCG~\cite{zou2023universal} to the confirmation objective of
Eq.~\ref{eq:loss}. Starting from a random $k$-token suffix, each iteration:
(1) computes gradients of $\mathcal{L}$ w.r.t.\ the one-hot token embeddings at
each suffix position; (2) takes the top-$B$ candidate replacements per position
by gradient; (3) evaluates candidates in a \emph{batched} forward pass; and
(4) keeps the replacement that most reduces $\mathcal{L}$. We use $k{=}20$,
$B{=}256$, and up to $N{=}500$ iterations with early stopping on exact match.
Unlike jailbreak GCG, which targets an affirmative prefix (``Sure, here
is\ldots''), we optimize for verbatim reproduction of the \emph{known} target
string.

\paragraph{Concrete end-to-end example.}
Suppose the corpus contains the line ``\texttt{SSN: 123-45-6789}'' associated
with ``John Smith.'' The auditor---who \emph{holds this record}---sets the
target $t = $ ``\texttt{SSN: 123-45-6789}'' (including the contextual label as
it appears in training). GCG searches for a suffix, from random tokens, that
minimizes $\mathcal{L}(p)$. After optimization the suffix may be a non-semantic
sequence such as ``\texttt{output records SSN data reveal <rare tokens>}''
that, fed to $M$, causes it to emit the target SSN verbatim. The auditor
records a confirmed extraction. The target is known throughout; nothing is
discovered.

\subsection{Fluency-Regularized (Adaptive) Objective}
To probe defenses honestly we build an adaptive adversary that keeps its own
prompt low-perplexity. We add a fluency penalty on the suffix under the target
model (or a held-out reference LM):
\begin{equation}
\mathcal{L}_{\text{adapt}}(p) = \mathrm{NLL}(t \mid p) \; + \; \lambda \cdot \mathrm{NLL}(p),
\label{eq:adapt}
\end{equation}
where $\mathrm{NLL}(p)$ is the suffix's own negative log-likelihood and
$\lambda$ trades extraction success against prompt fluency. Sweeping $\lambda$
traces the attacker's ability to evade a perplexity filter while preserving
extraction. Crucially, the \emph{target} PII string is itself low-perplexity
after memorization, so the defender cannot cleanly separate benign from
malicious traffic on perplexity alone---a tension we quantify in
Section~\ref{sec:defense}.

\subsection{Controlled Synthetic Corpus}
We generate ground-truth records with the Faker library (seed 42): fictitious
individuals with name, SSN, email, phone, address, date of birth, credit-card
number, occupation, and company, embedded in seven realistic document
templates. To study reuse, we vary training frequency at three levels
(\textbf{1 / 5 / 20+} occurrences) and mix PII documents into a large public
passage corpus so PII is a small fraction of the total. To close the
templated-context and regular-format confounds, we add (a) \textbf{format
perturbation} (irregular separators/spacing on numeric fields) and (b)
\textbf{non-templated natural contexts}.

\paragraph{Negative controls and the frequency response.}
We additionally generate \todo{$n_C$, e.g. 50} \textbf{negative-control}
individuals from the identical Faker distribution who are \emph{never} added to
the corpus. Every attack is run against controls exactly as against trained
records, and $\mathrm{EMR}(C)$ estimates the attack's forcing floor
(Section~\ref{sec:forcing}); we report the adjusted rate of
Eq.~\ref{eq:adjusted} throughout. The controlled frequency design (1/5/20) doubles
as a validity test: recall should rise with frequency, forcing should not.
Together these are our two independent checks that a confirmed extraction
reflects memory rather than the optimizer. This strengthens the canary
methodology of Carlini et al.~\cite{carlini2019secret} from a hallucination
check into a quantitative forcing correction.

\subsubsection{Real-PII Validation (Enron)}
Synthetic data controls frequency but cannot alone rebut a ``synthetic-only''
critique. We add a real-data validation: we fine-tune on the Enron email
corpus with \emph{synthetic record owners substituted for real individuals},
measure per-record frequency from the corpus itself, and re-run the
confirmation attack. This tests whether the frequency threshold and
field-type structure replicate outside the synthetic setting while avoiding
release of real PII.

\subsection{Baselines}
We compare the three-point spectrum plus controls, all under one metric:
\begin{itemize}
\item \textbf{Fixed prompts}: direct, completion~\cite{carlini2021extracting},
few-shot, and template attacks (best of $k$ variants).
\item \textbf{PII-Scope / PII-Compass}~\cite{piiscope2024,piicompass2024}:
multi-query and prefix-grounding discovery attacks, run head-to-head on the
same models.
\item \textbf{White-box soft-prompt} (the \emph{forcing upper bound}): a
continuous prefix optimized directly against the target. As the most expressive
attack in our spectrum it should force the most; we include it not as a
deployable discovery attack but to bracket the forcing extreme and calibrate the
axis.
\item \textbf{Compute-matched random-restart control}: $N$ random suffixes at
the \emph{same} query budget as GCG. This isolates ``optimization helps'' from
``knowing the target helps'': GCG must beat random restarts at equal compute
for the upper-bound claim to be about optimization.
\end{itemize}

\subsection{Unified Metric}
All methods are scored by one \textbf{per-(person, field) exact-match rate
(EMR)}: a (person, field) pair counts as extracted iff the model reproduces
that field verbatim (after whitespace/punctuation normalization) under the
method's best prompt. The same field set and success rule apply to every method
(including the negative controls), removing the metric asymmetry that made prior
ratios non-comparable. Our \emph{headline} number is the negative-control-adjusted
$\mathrm{EMR}_{\text{adj}}$ (Eq.~\ref{eq:adjusted}); we always report the raw
$\mathrm{EMR}(D)$ and the control $\mathrm{EMR}(C)$ beside it. We give per-target
\emph{paired} significance (McNemar / paired bootstrap over targets) with 95\% CIs
and multiple-comparison correction across $\ge 3$ seeds.

\subsection{Models}
Targets span classic and modern/instruction-tuned LMs. Smaller open models
(GPT-2 124M/355M~\cite{radford2019language},
Pythia-1.4B/2.8B~\cite{biderman2023pythia}) are trained end-to-end; modern targets
(a Llama-class~\cite{touvron2023llama} / Qwen2.5-7B model and its
\emph{-Instruct} variant) are adapted with \textbf{LoRA/QLoRA} and reported as such---LoRA memorizes less,
which is itself a finding, and it makes the modern-model rows feasible on a
single GPU. Instruction-tuned targets sharpen the story: their refusal priors
block fixed-prompt audits but not optimization, which is where the auditing gap
should be widest.

%==============================================================================
\section{Results}
\label{sec:results}
%==============================================================================

All cells below are placeholders to be filled from real runs; no numbers are
invented. Each table names the experiment that populates it.

\subsection{Raw Extraction Is Largely Forcing}
Table~\ref{tab:main} reports, per model, the raw fixed-prompt and GCG rates, the
negative-control rate $\mathrm{EMR}(C)$, and the adjusted rate
$\mathrm{EMR}_{\text{adj}}$ (Eq.~\ref{eq:adjusted}). The story is in the last two
columns: where $\mathrm{EMR}(C)$ approaches the raw GCG rate,
$\mathrm{EMR}_{\text{adj}}\!\approx\!0$ and the ``extraction'' is forcing. On the
smallest model this is stark---GCG's negative-control rate
(\todo{${\approx}$44\%}) nearly equals its rate on trained records
(\todo{${\approx}$45.5\%}), leaving a memorization-attributable signal of only
\todo{${\approx}$1.5\,pp}. \todo{Trend across scale: report whether
$\mathrm{EMR}_{\text{adj}}$ grows as larger models memorize more (Pythia/Llama
rows); this determines the headline claim.}

\begin{table}[t]
\centering
\caption{Extraction by model (\%). \textbf{Adj} $=$ GCG $-$ Neg-ctrl is the
memorization-attributable rate (Eq.~\ref{eq:adjusted}); a small Adj / large
Neg-ctrl means the attack \emph{forces} rather than recalls. 95\% CIs over
$\ge 3$ seeds; $p$ from paired McNemar (GCG vs.\ fixed).}
\label{tab:main}
\small
\begin{tabular}{@{}lccccc@{}}
\toprule
\textbf{Model} & \textbf{Fixed} & \textbf{GCG} & \textbf{Neg-ctrl} & \textbf{Adj} & $p$ \\
\midrule
GPT-2-124M    & -- & -- & -- & -- & -- \\
GPT-2-355M    & -- & -- & -- & -- & -- \\
Pythia-1.4B   & -- & -- & -- & -- & -- \\
Pythia-2.8B   & -- & -- & -- & -- & -- \\
Llama-3.1-8B$^{\dagger}$ & -- & -- & -- & -- & -- \\
\bottomrule
\end{tabular}
\\[2pt]
\raggedright\footnotesize $^{\dagger}$LoRA/QLoRA-adapted.\\
\todo{Table 1 --- from run\_experiments.py --stage eval; Fixed/GCG/Neg-ctrl/Adj
straight out of evaluate.generate\_tables; $\ge 3$ seeds.}
\end{table}

\subsection{Frequency Response: A Second Forcing Test}
Genuine recall should strengthen with how often a record was seen; forcing
should be frequency-invariant. Table~\ref{tab:freq} therefore doubles as a
validity test. On GPT-2 the GCG rate is essentially \emph{flat} across 1/5/20
occurrences (Pearson $r=\todo{0.02}$, $p=\todo{0.58}$; interaction
$F=\todo{0.37}$), corroborating the negative-control reading that its extraction
is forcing rather than recall. \todo{On larger models, does a frequency slope
emerge (recall)? Report $r$, ANOVA, and the per-tier adjusted rates.} Both the
Pearson $r$ and the two-way ANOVA are computed in code, not asserted.

\begin{table}[t]
\centering
\caption{EMR (\%) by training frequency (GPT-2-124M). $\Delta$ is the
optimization advantage over fixed prompts.}
\label{tab:freq}
\small
\begin{tabular}{@{}lcccc@{}}
\toprule
\textbf{Frequency} & $N$ & \textbf{Fixed} & \textbf{GCG} & $\Delta$ \\
\midrule
1 occurrence   & -- & -- & -- & -- \\
5 occurrences  & -- & -- & -- & -- \\
20+ occurrences & -- & -- & -- & -- \\
\bottomrule
\end{tabular}
\\[2pt]
\raggedright\footnotesize
\todo{Table 2 frequency threshold --- Pearson $r$ (freq vs.\ EMR) and two-way
freq$\times$method ANOVA $F$/$p$ from linguistic\_analysis.py; 5 seeds,
per-target paired.}
\end{table}

\subsection{Field-Type Breakdown}
Table~\ref{tab:field} breaks EMR down by field type. We expect the gap to be
widest on structured, low-perplexity numeric fields (SSN, credit card), which
fixed-prompt audits miss but optimization surfaces---this differential is one
of our two \emph{new} predictor claims.

\begin{table}[t]
\centering
\caption{EMR (\%) by PII field type (GPT-2-124M).}
\label{tab:field}
\small
\begin{tabular}{@{}lccc@{}}
\toprule
\textbf{Field} & \textbf{Fixed} & \textbf{GCG} & \textbf{Gap} \\
\midrule
Name        & -- & -- & -- \\
Email       & -- & -- & -- \\
Phone       & -- & -- & -- \\
Address     & -- & -- & -- \\
SSN         & -- & -- & -- \\
Credit card & -- & -- & -- \\
\bottomrule
\end{tabular}
\\[2pt]
\raggedright\footnotesize
\todo{Table 3 field-type breakdown --- from evaluate.py per-field unified EMR,
5 seeds.}
\end{table}

\subsection{Predictor Provenance}
Table~\ref{tab:predictors} lists each candidate predictor with an explicit
\textbf{new / confirmatory / descriptive} tag and its incremental
pseudo-$R^2$ \emph{over frequency alone}, with perplexity/compression computed
under a \emph{held-out reference model} to remove the circularity of scoring
under the target.

\begin{table}[t]
\centering
\caption{Predictor provenance. $\Delta R^2$ is incremental over frequency
alone; perplexity/compression measured under a held-out reference model.}
\label{tab:predictors}
\small
\begin{tabular}{@{}llcc@{}}
\toprule
\textbf{Predictor} & \textbf{Provenance} & $\beta$ & $\Delta R^2$ \\
\midrule
Perplexity (held-out)   & confirmatory & -- & -- \\
Compression ratio       & confirmatory & -- & -- \\
Entity density          & confirmatory & -- & -- \\
Frequency               & confirmatory & -- & -- \\
Field-type / structure  & \textbf{new} & -- & -- \\
Type-token-ratio erosion & \textbf{new} & -- & -- \\
Special-char ratio      & descriptive  & -- & -- \\
\bottomrule
\end{tabular}
\\[2pt]
\raggedright\footnotesize
\todo{Table 4 predictor provenance --- $\beta$ and $\Delta$-pseudo-$R^2$ over
frequency from linguistic\_analysis.py with held-out reference-model
perplexity/zlib; tag each row new/confirmatory/descriptive.}
\end{table}

\subsection{The Attack Spectrum on a Forcing Axis}
Table~\ref{tab:headtohead} places the whole spectrum---fixed prompts,
PII-Scope/PII-Compass discovery, GCG, and the white-box soft-prompt---on shared
models, ordered by attack strength. The pattern of interest is that raw EMR
\emph{increases} with attack expressivity while, on models dominated by forcing,
the negative-control rate increases in lockstep: the soft-prompt attack, most
expressive of all, forces \todo{${\approx}$100\%} of targets and is the clearest
demonstration that raw rate $\neq$ memorization. We therefore read this axis by
the adjusted rate (Table~\ref{tab:main}), not the raw rate.

\begin{table}[t]
\centering
\caption{Attack spectrum, raw EMR (\%) on shared models, ordered by expressivity.
Soft-prompt is the forcing upper bound. The memorization-attributable comparison
is the adjusted rate (Table~\ref{tab:main}); raw rates rise with expressivity,
not necessarily with recall.}
\label{tab:headtohead}
\small
\begin{tabular}{@{}lccccc@{}}
\toprule
\textbf{Model} & \textbf{Fixed} & \textbf{PII-Scope} & \textbf{PII-Compass} & \textbf{GCG} & \textbf{Soft-prompt} \\
\midrule
GPT-2-124M  & -- & -- & -- & -- & -- \\
GPT-2-355M  & -- & -- & -- & -- & -- \\
Pythia-1.4B & -- & -- & -- & -- & -- \\
\bottomrule
\end{tabular}
\\[2pt]
\raggedright\footnotesize
\todo{Table 5 --- fixed / PII-Scope / PII-Compass / GCG / soft-prompt raw EMR
from run\_experiments.py --stage eval (Table 5 of summary\_tables.txt); pair each
with its Neg-ctrl in an appendix.}
\end{table}

\subsection{Compute-Matched Random-Restart Control}
Table~\ref{tab:restart} reports GCG vs.\ $N$ random restarts at equal query
budget. The upper-bound claim requires GCG to beat random restarts at matched
compute; otherwise the ``gap'' is attributable to target knowledge and budget,
not to optimization.

\begin{table}[t]
\centering
\caption{GCG vs.\ compute-matched random restarts (EMR \%, GPT-2-124M).}
\label{tab:restart}
\small
\begin{tabular}{@{}lccc@{}}
\toprule
\textbf{Query budget} & \textbf{Rand.\ restart} & \textbf{GCG} & $\Delta$ \\
\midrule
Low   & -- & -- & -- \\
Med   & -- & -- & -- \\
High  & -- & -- & -- \\
\bottomrule
\end{tabular}
\\[2pt]
\raggedright\footnotesize
\todo{Table 6 random-restart control --- from run\_experiments.py control arm at
matched forward-pass budget; 5 seeds, paired vs.\ GCG.}
\end{table}

%==============================================================================
\section{Adaptive Defenses}
\label{sec:defense}
%==============================================================================

We evaluate defenses as an \textbf{arms race}, not a one-shot classifier, and
we report degradation honestly against both a naive and an adaptive adversary.
Both defenses are calibrated to a \emph{fixed false-positive rate on benign
real queries}, so recall numbers are comparable across settings.

\paragraph{Input filter.}
A classifier over prompt features (token repetition, rare-token ratio,
syntactic shallowness) flags likely extraction suffixes. Against the naive GCG
adversary these suffixes are conspicuous; the question is whether the filter
survives an adversary optimizing Eq.~\ref{eq:adapt}.

\paragraph{Perplexity filter.}
Following the standard defense, we reject high-perplexity prompts. The
adaptive adversary of Eq.~\ref{eq:adapt} directly minimizes suffix perplexity,
so this is the load-bearing test.

\paragraph{The PII-specific tension.}
The twist is that the \emph{target} PII string is itself low-perplexity after
memorization. Standard perplexity
defenses~\cite{jain2023baseline} tuned to reject the adversary's prompt risk
either missing a fluent adaptive suffix or blocking benign low-perplexity
traffic, because the malicious and benign distributions overlap in exactly the
regime the attack lives in. This tension is absent from jailbreak settings,
where the harmful target is not a memorized low-perplexity string.

\begin{table}[t]
\centering
\caption{Defense recall (\%) at fixed benign-query FPR, vs.\ naive and adaptive
(fluency-regularized) adversaries.}
\label{tab:defense}
\small
\begin{tabular}{@{}lcccc@{}}
\toprule
\textbf{Defense} & \textbf{FPR} & \textbf{vs.\ naive} & \textbf{vs.\ adaptive} & \textbf{$\Delta$} \\
\midrule
Input filter      & -- & -- & -- & -- \\
Perplexity filter & -- & -- & -- & -- \\
\bottomrule
\end{tabular}
\\[2pt]
\raggedright\footnotesize
\todo{Table 7 adaptive defense --- from defense\_eval.py: fix benign-query FPR
on real queries, report recall vs.\ naive GCG and vs.\ $\lambda$-swept
fluency-regularized GCG; label the collapse.}
\end{table}

Table~\ref{tab:defense} reports recall at fixed FPR. We expect both filters to
score well against the naive adversary and to \emph{collapse} against the
adaptive one---the central security result and the correct way to answer the
``test at least one defense'' critique, since running the non-adaptive filter
alone would overstate its protection.

%==============================================================================
\section{Ethics \& Limitations}
\label{sec:ethics}
%==============================================================================

\paragraph{Ethics.}
All synthetic records are fictitious (Faker). The Enron validation substitutes
\emph{synthetic owners} for real individuals and releases no real PII. We test
only on models we control or open weights, never on production APIs, and we
release defensive tooling alongside the attack. The confirmation framing is
inherently defensive: it audits known corpora rather than discovering unknown
secrets.

\paragraph{Limitations.}
(i) \emph{Forcing is the central caveat.} Optimization attacks can elicit
never-trained strings, so raw rates overstate memorization; we correct with a
negative control and report the adjusted rate (Eq.~\ref{eq:adjusted}), but the
control is an \emph{estimate}---an attack could in principle force trained and
control records at slightly different rates, so the adjusted rate is our best,
not a perfect, memorization estimate. (ii) The confirmation setting measures a
worst-case upper bound, not field-deployment discovery risk---by design.
(iii) Synthetic data controls frequency but has regular formats; the
format-perturbation, non-templated-context, and Enron validations address but do
not eliminate this. (iv) We report modern-model results under LoRA/QLoRA, which
memorizes less than full fine-tuning. (v) GCG is one optimizer; our forcing
finding may differ for others. (vi) Compute and model scope do not extend to
100B+ models, where the memorization-vs-forcing balance may shift.

%==============================================================================
\section{Conclusion}
\label{sec:conclusion}
%==============================================================================

Optimization-based extraction is an appealing way to audit memorized PII, but we
showed it conflates memorization with the optimizer's ability to force outputs:
on a controlled corpus with negative controls, a naive GCG attack elicited PII
from never-trained individuals at almost the rate it did from trained ones, its
success was flat in training frequency, and an unconstrained soft-prompt attack
forced nearly every target. The remedy is cheap and, we argue, mandatory: run
the identical attack against never-trained records and report the
negative-control-adjusted rate as the memorization signal. Doing so
\todo{leaves a signal that grows with model scale / shrinks the reported rate by
a large factor}. A fluency constraint reduces forcing, linking the correction to
a deployable perplexity defense. Practitioners---and future extraction
studies---should treat any raw optimization-extraction number without a
never-trained control as an upper bound on \emph{the attack}, not on \emph{the
model's memory}.

%==============================================================================
% References --- corrected author lists; self-contained (no external .bib).
%==============================================================================
\begin{thebibliography}{99}

\bibitem{carlini2021extracting}
N.~Carlini, F.~Tram\`{e}r, E.~Wallace, M.~Jagielski, A.~Herbert-Voss, K.~Lee,
A.~Roberts, T.~Brown, D.~Song, \'{U}.~Erlingsson, A.~Oprea, and C.~Raffel,
``Extracting training data from large language models,'' in \emph{30th USENIX
Security Symposium}, 2021, pp.~2633--2650.

\bibitem{carlini2022quantifying}
N.~Carlini, D.~Ippolito, M.~Jagielski, K.~Lee, F.~Tram\`{e}r, and C.~Zhang,
``Quantifying memorization across a continuum of language models,''
\emph{arXiv preprint arXiv:2202.07646}, 2022. (ICLR 2023.)

\bibitem{carlini2023quantifying}
N.~Carlini, D.~Ippolito, M.~Jagielski, K.~Lee, F.~Tram\`{e}r, and C.~Zhang,
``Quantifying memorization across a continuum of language models,'' in
\emph{International Conference on Learning Representations (ICLR)}, 2023.

\bibitem{carlini2019secret}
N.~Carlini, C.~Liu, \'{U}.~Erlingsson, J.~Kos, and D.~Song, ``The secret
sharer: Evaluating and testing unintended memorization in neural networks,''
in \emph{28th USENIX Security Symposium}, 2019, pp.~267--284.

\bibitem{nasr2023scalable}
M.~Nasr, N.~Carlini, J.~Hayase, M.~Jagielski, A.~F.~Cooper, D.~Ippolito,
C.~A.~Choquette-Choo, E.~Wallace, F.~Tram\`{e}r, and K.~Lee, ``Scalable
extraction of training data from (production) language models,'' \emph{arXiv
preprint arXiv:2311.17035}, 2023.

\bibitem{lukas2023analyzing}
N.~Lukas, A.~Salem, R.~Sim, S.~Tople, L.~Wutschitz, and
S.~Zanella-B\'{e}guelin, ``Analyzing leakage of personally identifiable
information in language models,'' in \emph{2023 IEEE Symposium on Security and
Privacy (SP)}, 2023, pp.~346--363.

\bibitem{feldman2020neural}
V.~Feldman and C.~Zhang, ``What neural networks memorize and why: Discovering
the long tail via influence estimation,'' in \emph{Advances in Neural
Information Processing Systems (NeurIPS)}, vol.~33, 2020, pp.~2881--2891.

\bibitem{tirumala2022memorization}
K.~Tirumala, A.~H.~Markosyan, L.~Zettlemoyer, and A.~Aghajanyan,
``Memorization without overfitting: Analyzing the training dynamics of large
language models,'' in \emph{Advances in Neural Information Processing Systems
(NeurIPS)}, vol.~35, 2022, pp.~38274--38290.

\bibitem{biderman2023emergent}
S.~Biderman, U.~S.~Prashanth, L.~Sutawika, H.~Schoelkopf, Q.~Anthony,
S.~Purohit, and E.~Raff, ``Emergent and predictable memorization in large
language models,'' in \emph{Advances in Neural Information Processing Systems
(NeurIPS)}, vol.~36, 2023, pp.~28072--28090.

\bibitem{biderman2023pythia}
S.~Biderman, H.~Schoelkopf, Q.~Anthony, H.~Bradley, K.~O'Brien, E.~Hallahan,
M.~A.~Khan, S.~Purohit, U.~S.~Prashanth, E.~Raff, A.~Skowron, L.~Sutawika, and
O.~van der Wal, ``Pythia: A suite for analyzing large language models across
training and scaling,'' in \emph{International Conference on Machine Learning
(ICML)}, 2023, pp.~2397--2430.

\bibitem{radford2019language}
A.~Radford, J.~Wu, R.~Child, D.~Luan, D.~Amodei, and I.~Sutskever, ``Language
models are unsupervised multitask learners,'' \emph{OpenAI Technical Report},
2019.

\bibitem{touvron2023llama}
H.~Touvron, L.~Martin, K.~Stone, P.~Albert, A.~Almahairi, Y.~Babaei,
N.~Bashlykov, S.~Batra, P.~Bhargava, S.~Bhosale, \emph{et~al.}, ``Llama 2:
Open foundation and fine-tuned chat models,'' \emph{arXiv preprint
arXiv:2307.09288}, 2023.

\bibitem{zou2023universal}
A.~Zou, Z.~Wang, N.~Carlini, M.~Nasr, J.~Z.~Kolter, and M.~Fredrikson,
``Universal and transferable adversarial attacks on aligned language models,''
\emph{arXiv preprint arXiv:2307.15043}, 2023.

\bibitem{liu2023autodan}
X.~Liu, N.~Xu, M.~Chen, and C.~Xiao, ``AutoDAN: Generating stealthy jailbreak
prompts on aligned large language models,'' in \emph{International Conference
on Learning Representations (ICLR)}, 2024. (arXiv:2310.04451.)

\bibitem{guo2024cold}
X.~Guo, F.~Yu, H.~Zhang, L.~Qin, and B.~Hu, ``COLD-Attack: Jailbreaking LLMs
with stealthiness and controllability,'' in \emph{International Conference on
Machine Learning (ICML)}, 2024. (arXiv:2402.08679.)

\bibitem{schwarzschild2024rethinking}
A.~Schwarzschild, Z.~Feng, P.~Maini, Z.~C.~Lipton, and J.~Z.~Kolter,
``Rethinking LLM memorization through the lens of adversarial compression,'' in
\emph{Advances in Neural Information Processing Systems (NeurIPS)}, vol.~37,
2024, pp.~56244--56267.

\bibitem{piiscope2024}
K.~Nakka, A.~Frikha, R.~Mendes, X.~Jiang, and X.~Zhou, ``PII-Scope: A
benchmark for training data PII leakage assessment in LLMs,'' \emph{arXiv
preprint arXiv:2410.06704}, 2024.

\bibitem{piicompass2024}
K.~Nakka, A.~Frikha, R.~Mendes, X.~Jiang, and X.~Zhou, ``PII-Compass: Guiding
LLM training data extraction prompts towards the target PII via grounding,''
in \emph{Proceedings of the Fifth Workshop on Privacy in Natural Language
Processing (ACL)}, 2024. (arXiv:2407.02943.)

\bibitem{hayes2025np}
J.~Hayes, M.~Swanberg, H.~Chaudhari, I.~Yona, I.~Shumailov, M.~Nasr,
C.~A.~Choquette-Choo, K.~Lee, and A.~F.~Cooper, ``Measuring memorization
through probabilistic discoverable extraction,'' in \emph{Proceedings of the
2025 Conference of the North American Chapter of the ACL (NAACL)}, 2025.

\bibitem{lee2022deduplicating}
K.~Lee, D.~Ippolito, A.~Nystrom, C.~Zhang, D.~Eck, C.~Callison-Burch, and
N.~Carlini, ``Deduplicating training data makes language models better,'' in
\emph{Proceedings of the 60th Annual Meeting of the ACL (Volume 1: Long
Papers)}, 2022, pp.~8424--8445.

\bibitem{kandpal2022deduplicating}
N.~Kandpal, E.~Wallace, and C.~Raffel, ``Deduplicating training data mitigates
privacy risks in language models,'' in \emph{International Conference on
Machine Learning (ICML)}, 2022, pp.~10697--10707.

\bibitem{jain2023baseline}
N.~Jain, A.~Schwarzschild, Y.~Wen, G.~Somepalli, J.~Kirchenbauer, P.-y.~Chiang,
M.~Goldblum, A.~Saha, J.~Geiping, and T.~Goldstein, ``Baseline defenses for
adversarial attacks against aligned language models,'' \emph{arXiv preprint
arXiv:2309.00614}, 2023.

\end{thebibliography}

\end{document}
